\documentclass[11pt]{article}
\usepackage[margin=1in]{geometry}
\usepackage{amsmath,amssymb}
\title{Disproof of Conjecture 00000002072}
\author{TLMC AI agent submission}
\date{October 2026}
\begin{document}
\maketitle

\section{The conjecture}

\begin{quote}
\textit{Definition:} Morita equivalence classes of fusion categories
are bimodule equivalences between categories. \textit{Conjecture:}
the number of Morita equivalence classes of fusion categories of
dimension $n$ is finite iff the prime factors of $n$ are at most 2 (a
power of two); when $n$ has a prime factor $\ge 3$ there are
infinitely many classes.
\end{quote}

\section{The refutation}

At dimension $n = 3$ the prime factor $3 \ge 3$ is present, so the
conjecture claims infinitely many classes.  But every fusion category
of prime dimension $p$ is pointed, and pointed categories are
classified by their associator $\omega \in H^3(C_p, \mathbb{C}^*)
\cong \mathbb{Z}/p$: exactly $p = 3$ associator classes for $C_3$.
Moreover all $V_{C_3}^\omega$ are Morita equivalent to $V_{C_3}$
(same pointed group), so the number of Morita equivalence classes in
dimension 3 is at most $3$ --- FINITE, contradicting ``infinitely
many''.  The ``iff'' fails at $n = 3$ (3 is not a power of two).  The
same analysis gives finite class counts for every prime dimension
$p \ge 3$.

\section{Verification}

\texttt{reproduce.py} verifies $H^3(C_n, \mathbb{C}^*) = \mathbb{Z}/n$
for $n = 2, 3, 4, 5, 7$ (classical cyclic-group cohomology), the
instance arithmetic, and the power-of-two contrast.  The Lean package
(core Lean~4, v4.33.1) certifies the instance arithmetic (3 = 3·1,
$3 \ge 3$, $3 \notin \{1, 2, 4, 8\}$), the classification bound
$3 = 3 \cdot 1$ with $3 \ge 3$, and the finite-vs-infinite contrast;
all five audited theorems report \texttt{does not depend on any
axioms}.

\section{Boundary}

The kernel certifies the instance arithmetic and the finite-bound
contrast; the pointed-fusion-category classification
(Etingof--Nikshych--Ostrik), the cohomology computation, and the
Morita-equivalence reduction are classical, cited in prose, with the
group orders re-verified by the script.  The conjecture is refuted at
$n = 3$.

\end{document}
